\section{Security Analysis}
\label{sec:security}

Property~\ref{prop:text} and Theorem~\ref{prop:loss} hold under A2, for a goal whose enrollment was authentic and whose signer is
uncompromised, with Inco satisfying I1 to I4. Property~\ref{prop:hold} is about the recorded deadline. Theorem~\ref{prop:leakproof} is about
the decision channel and states its own assumptions. The arguments are by hand over the transition model below; we have not mechanised them.
Section~\ref{sec:eval} reports how each is tested and with what.

\subsection{Transition model}
\label{sec:model}

Fix a goal opened with budget $B_0$, cap $C$, call count $n_0$, attempt budget $m_0$, allowlist $A$ and funding $F$. Each request $s$ is in
one state: \textsf{Pending}, \textsf{Rejected}, \textsf{Approved} (finalized, not funded), \textsf{Held} (funded, not captured),
\textsf{Paid}, or \textsf{Reclaimed}. A held payment the payer took back, or one past its deadline not yet voided, is still \textsf{Held} in
the vault's accounting until it is reclaimed (Figure~\ref{fig:states}). Write $a_s$ for its amount and define
\begin{itemize}
  \item $D = \sum_{s \in \textsf{Paid}} a_s$, the lifetime vendor payments;
  \item $R = \sum_{s \in \textsf{Approved} \cup \textsf{Held}} a_s$, the reserved, unresolved amounts;
  \item $Q = a_p \cdot [\,\mathit{ok}_p\,]$, the debit of the pending request $p$, if there is one;
  \item $L = |\textsf{Approved} \cup \textsf{Held} \cup \textsf{Paid}|$, the live approvals,
\end{itemize}
where $\mathit{ok}_p$ is the pending request's encrypted decision, which exists from the request on and is revealed later. Let $B$ be the
plaintext of the budget handle, $n$ the call count, $m$ the attempts left and $\sigma$ the last sequence number. The claims are
\begin{itemize}
  \item[(J1)] $B + D + R + Q = B_0$;
  \item[(J2)] $n + L = n_0$;
  \item[(J3)] $m + \sigma = m_0$;
  \item[(J4)] no request is both \textsf{Paid} and \textsf{Reclaimed}, and each is reclaimed at most once.
\end{itemize}
J2 counts \textsf{Paid} approvals and leaves the pending request out, because the call count moves at finalization.

\begin{table*}[t]
\centering\small
\begin{tabular}{@{}L{0.17\linewidth}L{0.31\linewidth}L{0.46\linewidth}@{}}
\toprule
Transition & Effect & Why J1 to J4 are preserved \\
\midrule
request $(a)$ & $\sigma{+}1$, $m{-}1$; $B \mathrel{-}= a\cdot\mathit{ok}$; new \textsf{Pending} & $Q$ rises by exactly what $B$ falls (J1); $m + \sigma$ fixed (J3). Only with no request pending, so $Q$ was 0. If $n = 0$, $\mathit{ok} = 0$ in plaintext. \\
finalize, approved & \textsf{Pending}$\to$\textsf{Approved}; $n{-}1$ & $Q$ moves into $R$ (J1); $L{+}1$, $n{-}1$ (J2). Needs $n \ge 1$, guaranteed since $\mathit{ok}$ required it at request (I1, I2). \\
finalize, rejected & \textsf{Pending}$\to$\textsf{Rejected} & $Q = 0$ before and after (J1). \\
authorize (fund) & \textsf{Approved}$\to$\textsf{Held} & Stays in $R$ and $L$. Refused unless approved, not reclaimed, goal open, terms and amount exact. \\
claim (capture) & \textsf{Held}$\to$\textsf{Paid} & $a$ moves from $R$ to $D$ (J1); stays in $L$ (J2). Escrow refuses after the capture deadline. \\
payer's reclaim & funds back to payer & No accounting change: still \textsf{Held} in the vault until credited. \\
reclaimExpired & \textsf{Approved}/\textsf{Held}$\to$\textsf{Reclaimed}; $B{+}a$, $n{+}1$ & $a$ moves from $R$ to $B$ (J1); $L{-}1$, $n{+}1$ (J2); $m$ unchanged (J3). Only in Table~\ref{tab:reclaim}'s cases, never from \textsf{Paid}; a flag stops a second credit (J4). \\
close, sweep & goal closed; payer's balance to owner & No accounting change; \code{authorize} and \code{requestEscrowSpend} then refuse. \\
\bottomrule
\end{tabular}
\caption{Every transition of a goal and why it preserves the invariants. A reclaim of an older approval while a newer request is
pending changes $R$ and $B$ only, and leaves $Q$ as it was.}
\label{tab:transitions}
\end{table*}

\begin{proposition}[Invariants]\label{prop:inv}
Under I1 to I4, J1 to J4 hold in every reachable state of a goal on the canonical chain.
\end{proposition}

\noindent\emph{Proof.} Initially $B = B_0$, $n = n_0$, $m = m_0$, $\sigma = 0$ and all sums are zero. Table~\ref{tab:transitions} lists
every function that changes a goal's state and shows each preserves J1 to J4; anything else reverts or leaves the goal unchanged.
Two points need the assumptions. J1 equates a plaintext with handle arithmetic, which needs I1 and I2: the budget handle each request reads
is the one the previous transaction wrote. And the \textsf{Approved}/\textsf{Rejected} split equals $\mathit{ok}$ only if the finalized value
is the attested one, which is I3 together with the vault verifying against the handle it stored. \hfill$\square$

\begin{property}[No authority from text]\label{prop:text}
No signature issued by the signer moves money to anyone other than the goal's owner unless the vault's record on the canonical chain holds
$(g, s)$ as finalized, approved and not reclaimed when the payment is funded, and the signed terms equal the terms recorded at request time.
\end{property}

\noindent The signer consumes no terms from its caller, and by Section~\ref{sec:binding} a signature moves tokens only inside the vault's
\code{authorize}, which re-derives and compares the whole payment, or \code{charge}, which pays only the owner of a closed goal.

\begin{theorem}[Lifetime loss bound]\label{prop:loss}
Under A2, for a goal whose enrollment was authentic and whose signer is uncompromised, with Inco satisfying I1 to I4, the lifetime vendor
payments satisfy $D \le \min(B_0, C\cdot n_0)$, every payment goes to a member of $A$, and each is at most $C$. Separately, vendor payments and
sweeps together never exceed the payer's funding $F$ plus any tokens others sent it: held funds that return to the payer are counted once, when
they are paid or swept. The bound holds whatever the signer's RPC reports and whatever blocks are dropped after the enrollment is stable.
Gas, the relay key's spending and lost availability are outside it.
\end{theorem}

\noindent\emph{Proof.} By J1, $D \le B_0$, since $B$, $R$ and $Q$ are non-negative. Every paid request was approved, and approval needs
$a \le C$ by I1 and I3, so each payment is at most $C$. Paid requests are live approvals, so $|\textsf{Paid}| \le L \le n_0$ by J2, and
$D \le C\cdot n_0$. Payees are checked against $A$ structurally at request, and capture pays the recorded payee. By Property~\ref{prop:text}
only records on the canonical chain count, which is why neither the RPC nor a dropped block can add to $D$. For the second bound, every
token that leaves the payer goes to the escrow or to the owner; the escrow pays out only by capture or returns to the payer, so the payer's
balance plus held funds plus $D$ plus sweeps equals $F$ plus inflows at all times. Authentic enrollment is needed because otherwise the payer
the user funds might not be the one the goal names; an uncompromised signer because the payer is a key. \hfill$\square$

\begin{property}[Recorded funding deadline]\label{prop:hold}
A signature for a request cannot fund the escrow at or after its recorded \textit{validBefore}, and the vendor cannot capture at or after the
recorded capture deadline. Whatever the orchestrator reports, no payment is funded an hour or more after its request, or captured two hours or
more after it.
\end{property}

\noindent The window is frozen in the record and is the one signed; the token and the escrow enforce it. This is a property of the recorded
deadline only. It protects a buyer against a late-collecting vendor if the orchestrator reports the hold honestly and with a margin for transit and
clocks; under A2 it weakens to the hour, without weakening Theorem~\ref{prop:loss}. Authenticated commercial holds need the vendor's signature on
an absolute deadline, as in Section~\ref{sec:delivery}.

\begin{property}[Reclaim cannot double-pay]\label{prop:reclaim}
For every $(g, s)$, the amount is never both captured and credited back, and is credited at most once.
\end{property}

\noindent This is J4. A credit needs nothing captured. If nothing was funded, it needs $t \ge \mathit{validBefore}$, after which funding is
impossible. If funds are held, it needs the capture deadline to have passed, after which capture is impossible, and the funds are voided in the
same transaction.

\subsection{The decision channel}
\label{sec:leakthm}

\begin{definition}[Decision-channel leakage]\label{def:leak}
The secret is $\mathsf{s} = (B_0, C)$ and $P$ is the public setup: everything public before the first request, including the funded amount,
expiry and allowlist. The pair $(\mathsf{s}, P)$ has any joint distribution, and every goal has attempt budget at most $m$. Before request $i$ an
adaptive adversary chooses which earlier approvals $\rho_i$ to reclaim and an amount $a_i$, as functions of $P$, its coins $r$ and everything
public so far, and observes $d_i = [\,n_i \ge 1\,] \wedge [\,a_i \le B_i\,] \wedge [\,a_i \le C\,]$. Its view after $k \le m$ requests is
$V = (\rho_1, a_1, d_1, \ldots, \rho_k, a_k, d_k)$. The leakage is $\mathcal{L} = H_\infty(\mathsf{s} \mid P) - H_\infty(\mathsf{s} \mid P, V)$, with
$H_\infty(X \mid Y) = -\log_2 \sum_y \max_x \Pr[x, y]$ the average conditional min-entropy~\cite{smith2009}.
\end{definition}

\begin{theorem}[Decision-channel bound]\label{prop:leakproof}
Assume that the adversary's coins are independent of $\mathsf{s}$ given $P$, and that no decision is revealed for a request that does not end on
the canonical chain. Then for every prior and every adversary strategy, $\mathcal{L} \le m$ bits.
\end{theorem}

\noindent\emph{Proof.} For each fixed $(p, r)$ each $a_i$, $\rho_i$ and the decision to stop are functions of $d_1, \ldots, d_{i-1}$, so the view is
determined by the decision bits. These form the leaves of a binary tree of depth at most $m$, so the view takes at most $2^m$ values $v$. Hence
$\sum_v \max_{\mathsf{s}} \Pr[\mathsf{s}, v \mid p, r] \le 2^m \max_{\mathsf{s}} \Pr[\mathsf{s} \mid p]$. Multiplying by $\Pr[p, r]$ and summing gives
$H_\infty(\mathsf{s} \mid P, V, r) \ge H_\infty(\mathsf{s} \mid P) - m$, the bound for deterministic channels of Smith~\cite{smith2009}. The second
assumption is what makes the depth $m$: an attempt is counted on the canonical chain, so a decision revealed for a request on a branch that is later
dropped would be an answer the count never charged. \hfill$\square$

\paragraph{From the channel to the deployed transcript.} The theorem bounds what the decisions reveal. Relating it to everything the deployed
system publishes needs further assumptions, which we state and do not prove. (i) Ciphertexts, handles, Inco's operation events and attestations
reveal nothing about $\mathsf{s}$ beyond the decisions; this is I5 together with handles being computed from operation inputs and not from
plaintexts, which is how Lightning derives them, and needs a computational argument we have not made. (ii) There are no timing or gas channels.
In one measurement on version~1's design both outcomes used 326{,}528 gas, and over 50 calls decision latency showed overlapping distributions,
which does not rule a channel out. (iii) No revealed decision is on a dropped branch; we have not measured how soon Inco reveals relative to
finality on Base. (iv) The user, client and console reveal nothing derived from the limits; Section~\ref{sec:otherleak} gives two ways an earlier
console did. (v) Each goal has its own secret; a user who reuses limits across goals gives an adversary up to the sum of their attempt budgets.

\paragraph{The bound is an average and a ceiling.} With a cap uniform on $0, \ldots, 1023$, a request for one unit that is rejected identifies
$C = 0$ from a single decision. And the bound can be close to tight: binary search narrows the cap's range by a factor of two per answer. Our
regression test, with a 10 USDC range, leaves an interval of 0.117 to 0.156 USDC after $m = 8$ and recovers the cap at $m = 30$.
Section~\ref{sec:leakeval} asks what the bound leaves in practice.

\subsection{What the bound does not cover}
\label{sec:otherleak}

\paragraph{Spending.} The vault never sees the 402 body, so a compromised orchestrator may request more than the seller asked, up to the cap, to an
allowlisted payee. A vendor that is paid and delivers nothing is paid. An allowlisted vendor that signs an inflated price is paid that price within
the limits. Removing the first needs a vendor-signed quote the vault verifies; the others are Section~\ref{sec:delivery}'s subject.

\paragraph{Leakage outside the decisions.} The payer's token balance is public, so funding it with the budget plus a margin, as an earlier console
did, publishes the budget before any request. The console now funds a fixed 5.00 USDC, above the largest budget it accepts, at the cost of a looser
fallback bound. An earlier console also priced its overcharge test vendor from the budget and passed that through the orchestrator. Every value a
deployment sends must be checked; the console's own defaults are constants in public code and so are not secrets.

\paragraph{Liveness.} A decision Inco never reveals blocks the goal, since only finalization clears a pending request; no money is at risk and the
owner can close the goal and sweep. On Base Sepolia, reveals need both of two covalidators, and the orchestrator refuses to commit a debit unless Inco
can reveal a known handle. A lying RPC can stall a run or waste a signature.
